% Sets and other symbols
\def\R{{\Bbb R}}
\def\N{{\Bbb N}}
\def\Z{{\Bbb Z}}
\def\Q{{\Bbb Q}}
\def\T{{\Bbb T}}
\def\C{{\Bbb C}}
\def\E{{\Bbb E}}
%\let\emptyset\varnothing % Bye-bye ugly zero
%\let\epsilon\varepsilon % Bye-bye not-very-mathy epsilon
\ifundefined{varsubsetneq}\let\varsubsetneq\subsetneq\fi
%\def\neq{\not\kern-.25pt=} % Fonts should be fixed now?
\def\neq{%
  \mathchoice
  {\kern.4pt\not\kern-.4pt=}
  {\kern.4pt\not\kern-.4pt=}
  {\kern1pt\not\kern-1pt=}
  {\kern1pt\not\kern-1pt=}}

% Functions
\def\mathfunc#1{\expandafter\def\csname #1\endcsname{\mathop{\rm\relax#1}}}
\mathfunc{Adj}
\mathfunc{Det}
\mathfunc{Aut}
\mathfunc{Ker}
\mathfunc{Dom}
\mathfunc{Gal}
\mathfunc{Img}
\mathfunc{Sup}
\mathfunc{Inf}
\mathfunc{Ord}
\mathfunc{Sgn}
\mathfunc{Ob}
\mathfunc{Irr}
\mathfunc{Deg}
\def\fractional{\mathop{\rm Frac}}
\def\pwr{{\scr P}}

% The following is useful for writing out units
% Turns everything into upright roman and turns `.' into cdot,
% but with appropriate spacing
\def\units#1{{
        \;
        \mathcode`.="2201
        \rm\relax#1
        }}

% Symbols / operators
\def\inv{^{-1}}
\def\of{\circ}
\def\degrees{^\circ}
\def\implies{\Longrightarrow}
\def\defeq{\mathrel{\buildrel\Delta\over=}}
\def\gline#1{{\buildrel \longleftrightarrow \over {#1}}}
%\def\qed{\hfill$\square$}
\def\qed{\hfill\vbox{\hrule width .5em height .5em depth 0em}}

\def\root #1 of #2{\hskip.3ex\raise.4ex\hbox{${}^{#1}$}\kern-1.4ex\sqrt{#2}}

\let\oldfrac\frac
\def\frac #1/#2{\hbox{\oldfrac #1/#2}}

% Calculus
\def\diff{{\rm d}}
\def\d/d#1{{\diff\over\diff#1}}
\def\integral from #1 to #2 of #3 d#4{%
\int_{#1}^{#2} #3 \,\diff #4%
}


% % Sectioning

% \def\section#1{%
%   \vskip\parskip
%   \vskip\parskip
%   \line{\sectionfont#1\hfill}
% }

% \def\newpage{\relax\vfill\eject}

% % Little sections
% \def\inset#1{{\ss\relax#1}}
% \newskip\insetleftskip
% \newskip\insetrightskip
% \newdimen\insetsize
% \insetleftskip=3em
% \insetrightskip=3em
% \insetsize=4.4in
% \everydisplay{\hskip\leftskip}

% \long\def\inset#1#2\endinset{%
%   \begingroup\par
%   \leftskip=\insetleftskip
%   \noindent\llap{\hbox to \insetleftskip{{\ss #1}\hfil}}#2\par
%   \endgroup}

% Unnumbered
\long\def\nonumdefinition#1{\inset{Definition.\hskip1em}{\it\relax#1}\endinset}
\long\def\nonumproposition#1{\inset{Proposition.\hskip1em}{\it\everymath={\/}\relax#1}\endinset}
\long\def\nonumtheorem#1{\inset{Theorem.\hskip1em}{\it\everymath={\/}\relax#1}\endinset}
\long\def\nonumconjecture#1{\inset{Conjecture.\hskip1em}{\it\everymath={\/}\relax#1}\endinset}
\long\def\nonumlemma#1{\inset{Lemma.\hskip1em}{\it\relax#1}\endinset}
\def\nonumquestion#1{\inset{Question.\hskip1em}{\it\everymath={\/}\relax#1}\endinset}
\def\nonumformula#1#2\endformula{\inset{{{\bf #1.}\hskip1em}}{\it\everymath={\/}\relax#2}\endinset}

\def\definitionword{Definition}
\def\propositionword{Proposition}
\def\theoremword{Theorem}
\def\questionword{Question}
\def\lemmaword{Lemma}
\def\formulaword{Formula}
\def\problemword{Problem}
% Numbered
\long\def\problem#1{\inset{Problem \subsubnums.\hskip1em}{\relax#1}\endinset}
\long\def\subproblem(#1)#2{\inset{(#1)\hskip1em}{\relax#2}\endinset}
\long\def\definition#1{\inset{Definition \subsubnums.\hskip1em}{\it\relax#1}\endinset}
\long\def\proposition#1{\inset{Proposition \subsubnums.\hskip1em}{\it\everymath={\/}\relax#1}\endinset}
\long\def\theorem#1{\inset{Theorem \subsubnums.\hskip1em}{\it\everymath={\/}\relax#1}\endinset}
\long\def\conjecture#1{\inset{Conjecture \subsubnums.\hskip1em}{\it\everymath={\/}\relax#1}\endinset}
\long\def\lemma#1{\inset{Lemma \subsubnums.\hskip1em}{\it\relax#1}\endinset}
\long\def\corollary#1{\inset{Corollary \subsubnums.\hskip1em}{\it\relax#1}\endinset}
\def\question#1{{\insetleftskip=8em\inset{\hbox to \insetleftskip{Question \subsubnums.\hfil}}{\it\everymath={\/}\relax#1}\endinset}}
\def\formula#1#2\endformula{\inset{{Formula \subsubnums. {\bf #1.}\hskip1em}}{\it\everymath={\/}\relax#2}\endinset}

% Neither numbered nor unnumbered
\long\def\axiom(#1)#2{{\insetleftskip=4em\inset{\hbox to \insetleftskip{\hfil\bf(#1)\hskip.5em}}{\it\relax#2}\endinset}}
\long\def\note#1{{\inset{Note.\hskip1em}{\relax#1}\endinset}}
\long\def\remark#1{{\inset{Remark.\hskip1em}{\relax#1}\endinset}}
\long\def\proof#1\endproof{{\insetleftskip=4em\inset{\hbox to \insetleftskip{Proof.\hfill}}#1\qed\endinset}}


% And for their numbered variants.  We count definitions,
% propositions, theorems, examples, and lemmas all using the same
% counter, so that e.g., 4.5.1 refers to exactly one thing and is
% not ambiguous.  In all of these, of the section or subsection
% number is zero then we are not in a section or subsection so we
% omit those numbers in the numbering.

\def\margin#1{%
\itemitem{\ss\relax#1}
}

\def\case#1{[{\it\relax#1}]}
\newbox\casebox
\def\newcase#1{%
  \setbox\casebox=\hbox to 5em plus 1fill{[{\it\relax#1}]}
  \box\casebox
}

\let\oldeqno\eqno
\def\eqno#1{\oldeqno{#1\relax\hskip\the\rightskip}}

